% =====================================================================
%  Short paper (max. 5 pages) -- IEEE conference format (IEEEtran)
%  PPI Grupo 4 -- UPeU Juliaca -- PeMS-D4 anomaly detection protocol
%  Version of 27 Oct 2026 (end of week 8 of the 12-week plan).
%
%  STATUS LEGEND used in the text:
%    "obtained"  -> results produced and verified by automatic assert
%                   checks (PPI Annex C, weeks 1--8).
%    "pending"   -> analyses defined a priori but not yet executed.
%  Lines marked  %% TODO  require information the authors must supply.
% =====================================================================
\documentclass[conference]{IEEEtran}
% \IEEEoverridecommandlockouts  % uncomment only if a funding \thanks is added
\usepackage{cite}
\usepackage{amsmath,amssymb,amsfonts}
\usepackage{graphicx}
\usepackage{textcomp}
\usepackage{xcolor}
\usepackage{booktabs}
\usepackage{pifont}

\newcommand{\cm}{\ding{51}}          % reported / satisfied
\newcommand{\xm}{\ding{55}}          % not done
\newcommand{\nr}{n.r.}               % not reported
\newcommand{\pend}{\textit{pending}} % analysis not yet executed
\newcommand{\pl}{\textit{plan}}         % planned, not yet executed
\usepackage{url}

\begin{document}

\title{Deep Learning versus a Non-Parametric Baseline for Freeway Traffic-Flow
Anomaly Detection: A Uniform Protocol and Labeling-Sensitivity Design on PeMS-D4}

%% TODO: confirm author order, e-mails or ORCID, and whether the advisor is
%% listed as co-author according to the conference rules.
\author{\IEEEauthorblockN{Wilbert Alex Mayta Arotaype}
\IEEEauthorblockA{\textit{Escuela Profesional de Ingenier\'ia de Sistemas} \\
\textit{Universidad Peruana Uni\'on}\\
Juliaca, Peru \\
email address or ORCID}
\and
\IEEEauthorblockN{Juan Enders Prieto Quispe}
\IEEEauthorblockA{\textit{Escuela Profesional de Ingenier\'ia de Sistemas} \\
\textit{Universidad Peruana Uni\'on}\\
Juliaca, Peru \\
email address or ORCID}
\and
\IEEEauthorblockN{Wrangler Choquechambi Luque}
\IEEEauthorblockA{\textit{Escuela Profesional de Ingenier\'ia de Sistemas} \\
\textit{Universidad Peruana Uni\'on}\\
Juliaca, Peru \\
email address or ORCID}
\and
\IEEEauthorblockN{Esteban Tocto Cano}
\IEEEauthorblockA{\textit{Escuela Profesional de Ingenier\'ia de Sistemas} \\
\textit{Universidad Peruana Uni\'on}\\
Juliaca, Peru \\
email address or ORCID}
}

\maketitle

% NOTE (template rule): no symbols, special characters, footnotes or math
% in the title or abstract.
\begin{abstract}
Deep learning detectors for traffic anomalies are usually reported under
heterogeneous datasets, splits, thresholds and labels, while recent
benchmark studies show that simple statistical detectors are often
competitive once evaluation flaws are removed. We present a uniform
protocol, whose main rules were fixed before any network result was seen,
that compares six
deep learning families (LSTM, GRU, CNN-1D, autoencoder, Transformer and a
CNN-LSTM-GRU hybrid) and one published model (SI-DMSeq) against a
non-parametric rolling-statistics baseline on PeMS-D4 freeway flow series.
The protocol uses walk-forward temporal validation, a common reduced tuning
procedure for the six trained families with the published model left at its
default configuration, three seeds for the final models, PR-AUC without
point adjustment, paired tests with Holm correction, latency benchmarking,
and two separately constructed labeling schemes, a causal statistical proxy
and the controlled injection of point, contextual and collective anomalies,
to test whether the detector ranking depends on how anomalies are defined.
Preliminary outputs, which passed the specified automated consistency
checks, cover data curation (291 of 307 sensors retained), both labeling
schemes and the baseline. The proxy scheme yields a prevalence of 0.10
percent and only 29.2 percent of evaluation blocks contain both classes; its
positive set barely overlaps with the injected one (Jaccard index of at most
0.0012); and the baseline reaches a block-mean test PR-AUC of 0.345 under
injection, about 11 times chance, but a pooled PR-AUC of only 0.003 under the
proxy, about three times chance. Measured training times forced a reduced,
pre-declared hyperparameter search. Deep learning results and statistical
tests are pending.
\end{abstract}

\begin{IEEEkeywords}
anomaly detection, deep learning, evaluation protocol, PeMS, time series,
traffic flow, synthetic anomaly injection
\end{IEEEkeywords}

% ---------------------------------------------------------------------
\section{Introduction}

\textbf{Context and who is affected.} Traffic management centers, such as the
California Department of Transportation (Caltrans), which sponsors PeMS
\cite{chen2003pems}, rely on freeway sensors to detect disruptions, and
drivers bear the delays. The U.S.\ Federal Highway Administration attributes
about half of congestion to temporary disruptions, with incidents alone
accounting for 25\% of congestion \cite{fhwa_nonrecurring}. In the San
Francisco Bay Area, INRIX estimated that congestion cost each driver 79 hours
and USD~2\,250 in 2017, the year before the period studied here
\cite{cbs2018inrix}. Detecting atypical observations promptly therefore has
operational value. If detectors are adopted on the strength of results that
cannot be compared, an operator may pay for costly deep models without a
demonstrated advantage over a simple rule, which is why a uniform evaluation
protocol is needed.

Freeway loop detectors produce dense, seasonal and noisy flow series in
which timely identification of atypical observations supports traffic
monitoring. DL detectors based on recurrence, convolution, reconstruction or
attention have been proposed for this task
\cite{zhang2022automatic,mendez2022using,ijaradar2025datadriven,nouri2025sidmseq},
but their results are hard to compare: studies differ in data, splits,
labels, thresholds and metrics, and large public freeway datasets such as
PeMS were built for forecasting and carry no anomaly labels
\cite{coursey2024ftaed}. Labels are therefore derived from incident logs,
rules or simulation \cite{sun2022datacentric,yang2025sumo}, and a statistical
anomaly is not necessarily a real incident.

Two lines of evidence motivate a stricter evaluation. Point adjustment (PA)
inflates scores so much that random scores reach state-of-the-art F1
\cite{kim2022rigorous}, and popular benchmarks are largely solvable by
one-line moving-statistics rules \cite{wu2022flawed}. When evaluation is
controlled, classical detectors match or outperform deep ones on generic
benchmarks
\cite{schmidl2022anomaly,audibert2022deep,rewicki2023worth,liu2024elephant}
and on expert-labeled traffic volume \cite{taheri2026applicability};
moreover, label quality changes conclusions in traffic incident detection
\cite{sun2022datacentric}, and rankings shift with anomaly type and metric
\cite{liu2024elephant,sorbo2024navigating}.

\textbf{Research object and problem.} The object of study is DL-based
anomaly detection in freeway traffic flow, analyzed on PeMS-D4 with the
detector $\times$ labeling scheme $\times$ sensor-fold block as unit of
analysis. The problem is whether DL detectors provide a reproducible
advantage over a rolling-statistics baseline, and whether the resulting
ranking remains stable across operational definitions of anomaly and under
an inference-cost constraint.

\textbf{Knowledge gap.} Uniform comparisons already exist for generic series
\cite{schmidl2022anomaly,liu2024elephant} and for urban traffic volume
\cite{taheri2026applicability}. What is not known is whether the relative
standing of DL detectors against a simple baseline on freeway flow survives a
change in how anomalies are defined. In the literature we reviewed, no study
examines this on one temporal partition of a freeway dataset while jointly
comparing generic DL families, a non-parametric baseline and a published
traffic model, treating inference latency as a first-order criterion, and
quantifying seed variability with multiple-comparison control
(Table~\ref{tab:practices}). The gap is therefore one of evidence on
robustness, not of detectors.

\textbf{Objectives and questions.} \emph{General objective:} to determine
whether DL detectors provide a reproducible advantage over a rolling-statistics
baseline (G0) on PeMS-D4 and whether the detector ranking is stable across
labeling schemes. \emph{Specific objectives:} (O1)~estimate, under Scheme B,
the difference in PR-AUC between each DL detector and G0; (O2)~establish
which detectors that exceed G0 meet a predefined latency criterion;
(O3)~quantify the agreement between the detector rankings under Schemes A
and B. \emph{Research questions:} RQ1~(O1, H1), by how much, if at all, does
any DL detector improve average precision over G0 under Scheme B; RQ2~(O2,
H2), do the detectors that improve on G0 meet the latency criterion; RQ3~(O3,
H3), is the detector ranking stable across labeling schemes? Input-controlled
and per-type analyses test the robustness of these comparisons.

\textbf{Declared contributions.} (1)~\emph{Empirical:} evidence on whether
the standing of DL detectors against a rolling baseline on freeway flow is
robust to the anomaly definition (DL results pending); (2)~\emph{methodological:}
a reproducible protocol for PeMS-D4 that operationalizes practices distilled
from the seed studies, with two separately constructed labeling schemes and
an a~priori ranking-stability criterion; (3)~\emph{practical:} latency as an
explicit criterion, with a PR-AUC$\times$latency Pareto front. The
preliminary outputs reported here concern the data, the labels and the
baseline only.

% ---------------------------------------------------------------------
\section{Related Work and Distilled Practices}

\textbf{Traffic anomaly detection.} Reported detectors include
graph-attention forecasting \cite{zhang2022automatic}, CNN and BiLSTM
autoencoders \cite{mendez2022using}, a GCN-LSTM-attention autoencoder
evaluated with artificial anomalies \cite{ijaradar2025datadriven}, and
SI-DMSeq, a multitask LSTM encoder--decoder that flags disruptions with a
control chart \cite{nouri2025sidmseq}. Labeling is the main obstacle: weak
supervision raised the incident-class F1 of an LSTM from 0.69 to 0.88 on
PeMS District~4 \cite{sun2022datacentric}, and FT-AED combines crash logs
with expert labels \cite{coursey2024ftaed}. On expert-labeled volume data, a
mean$\pm c\cdot$std time-of-day profile reached F1 0.566 versus 0.120 for an
LSTM \cite{taheri2026applicability}.

\textbf{Evaluation pitfalls.} Under PA, uniform random scores reach
F1 0.969 on SWaT versus 0.216 without it \cite{kim2022rigorous}; 316 of 367
Yahoo series are separable by a moving mean/std one-liner
\cite{wu2022flawed}. AUC-ROC is dominated by true negatives under imbalance
\cite{sorbo2024navigating,liu2024elephant}, but point-wise PR-AUC is not
robust to small temporal lags \cite{liu2024elephant}.

\textbf{Synthetic labels.} Anomalies split into point-wise (global,
contextual) and pattern-wise types \cite{lai2021revisiting}. Injected
anomalies help to train and select detectors, but differ from real ones and
transfer only when types match \cite{carmona2022ncad,goswami2023unsupervised}.

Table~\ref{tab:practices} summarizes the protocol features reported by the
seed studies and motivates each design choice in Section~III.

\begin{table}[t]
\centering
\caption{Protocol features reported by seed studies.
Tr: traffic data; NB: non-deep baseline; CS: chronological split stated;
NT: no PA and no test-tuned threshold for the headline metric; RR: repeated
runs; ST: significance test; LT: runtime/latency.
\cm\ yes, \xm\ no, p partial, \nr\ not reported, -- not applicable,
plan: planned, not yet executed.}
\label{tab:practices}
\setlength{\tabcolsep}{3.2pt}
\renewcommand{\arraystretch}{0.9}
\footnotesize
\begin{tabular}{@{}lccccccc@{}}
\toprule
Study & Tr & NB & CS & NT & RR & ST & LT \\
\midrule
Taheri'26 \cite{taheri2026applicability} & \cm & \cm & \nr & \xm & \nr & \xm & \xm \\
Taheri'24 \cite{taheri2024tpfds}         & \cm & \cm & \nr & \xm & \cm & \xm & \xm \\
FT-AED'24 \cite{coursey2024ftaed}        & \cm & p   & \nr & \cm & \xm & \xm & \xm \\
Sun'22 \cite{sun2022datacentric}         & \cm & \xm & \cm & \cm & \cm & \xm & \xm \\
Zhang'22 \cite{zhang2022automatic}       & \cm & \cm & \nr & \cm & \xm & \xm & \xm \\
Ijaradar'25 \cite{ijaradar2025datadriven}& \cm & \cm & \nr & \nr & \nr & \xm & \xm \\
Nouri'25 \cite{nouri2025sidmseq}         & \cm & \xm & \cm & --  & \xm & \xm & \cm \\
Kim'22 \cite{kim2022rigorous}            & \xm & \cm & --  & \xm & p   & \xm & \xm \\
Audibert'22 \cite{audibert2022deep}      & \xm & \cm & --  & \xm & \nr & \cm & \xm \\
Schmidl'22 \cite{schmidl2022anomaly}     & \xm & \cm & --  & \cm & \nr & \xm & \cm \\
Liu'24 \cite{liu2024elephant}            & \xm & \cm & --  & \cm & \xm & \cm & \nr \\
\midrule
\textbf{This protocol}                   & \cm & \cm & \cm & \cm & \pl & \pl & \pl \\
\bottomrule
\end{tabular}
\end{table}

% ---------------------------------------------------------------------
\section{Protocol}

\textbf{Design.} Computational comparative experiment with a factorial
structure (detector $\times$ labeling scheme) and repeated measures over
sensor$\times$fold blocks, under walk-forward temporal validation and one
held-out test block.

\textbf{Data.} PeMS-D4 \cite{guo2019astgcn} contains flow, occupancy and
speed from 307 freeway sensors in the San Francisco Bay Area, 1~Jan--28~Feb
2018, at 5-min resolution (16\,992 steps per sensor). Sensors with more than
20\% zeros or a 24-h frozen run were excluded \emph{a priori}, leaving 291.
%% TODO: confirm that this quality control examined the full period before
%% the partition; if so, say it here (a reviewer asked for the declaration).
Each sample is a 12-step window ($t-11..t$, one hour) with $F=8$ inputs
(flow, occupancy, speed, cyclic hour and weekday encodings, a peak-hour
flag); the label is that of the final step $t$, so all methods score the same
instant. One global model per architecture is trained on the windows of all
sensors.

\textbf{Partition.} The last 20\% of each series is a held-out test block
used once. The first 80\% forms three expanding walk-forward folds; z-score
normalization, hyperparameters and thresholds are fitted only on the
training and validation parts of each fold, never on test data
(Table~\ref{tab:practices}, CS/NT).

\textbf{Labeling schemes.} \emph{Scheme A} (proxy) labels $x_{s,t}$ as
anomalous when $|z_{s,t}|>3$, with mean and standard deviation computed
causally on the previous 288 steps of raw flow; the anticipated prevalence was
0.3--2\%. The baseline G0 belongs to
the same rule family (rolling mean and standard deviation, although over a
much shorter window), so A is used only for the ranking-stability test (H3),
never for H1. \emph{Scheme B} injects point, contextual and collective
(3, 6 or 12 steps) anomalies into real series, perturbing flow, occupancy and
speed, at a 3\% target prevalence per split and per sensor$\times$fold block,
with a fixed injection seed, physical constraints (flow, speed $\geq 0$;
occupancy in $[0,1]$) and no overlap across partitions. Contextual anomalies
replace a value with one below the 5th or above the 95th percentile of a
per-sensor hour$\times$weekday profile estimated on training data, a rule
family distinct from G0 \cite{lai2021revisiting}.

\textbf{Conditions.} G0: non-parametric rolling detector on flow only; it
scores each step by the absolute z-score over the previous $w$ steps (window
$w\in\{6,11,22,44\}$ steps and threshold multiplier $k\in\{2.0,2.5,3.0\}$,
a 12-point grid). G1: LSTM; G2:
GRU; G3: CNN-1D; G4: autoencoder (reconstruction error as score); G5:
Transformer encoder; G6: CNN-LSTM-GRU hybrid; G7: SI-DMSeq
\cite{nouri2025sidmseq} with its published configuration, using the local
reconstruction error of the series at step $t$ as the point score. G1--G3, G5
and G6 optimize a supervised classification loss (sigmoid output, binary
cross-entropy). G4 optimizes a reconstruction loss on windows whose final-step
label is normal under the corresponding scheme, so those windows may contain
earlier anomalous steps. G0 and G7 optimize no supervised classification
loss, although label-based selection is used where specified (G0's window and
threshold on validation). Architecture and degree of label use are therefore
partly confounded, and the comparison is between detector systems rather
than between architectures in isolation. G0--G6 are implemented in
TensorFlow/Keras. Each of G1--G6 evaluates 12
configurations per scheme, the first being its initial specification and the
other 11 drawn from the fixed ranges of Table~\ref{tab:ranges} with a
deterministic seed. G7 keeps PyTorch
and no tuning, an asymmetry declared a~priori.

\textbf{Reduced search.} The nominal search (6 architectures $\times$ 2
schemes $\times$ 3 folds $\times$ 12 configurations $\times$ 3 seeds) is 1\,296
trainings; at the 4.3--5.4 min per epoch measured on fold 3 (CPU, LSTM), 30
epochs each exceed 70 days. Before any search we fixed one reduction for all
architectures: one seed, at most 10 epochs with early stopping (patience 3)
on PR-AUC, every fourth training window, and selection on the \emph{full}
validation block. The winner is retrained on the whole 80\% block with three
seeds, all windows and no early stopping, using the median best epoch. Cost was timed on an AMD Ryzen 7 7445HS
(6 cores, 12 threads) with 15.2\,GB of RAM, on CPU (TensorFlow 2.16.2); a
GPU run (RTX 4050, under WSL2) was slower for this model and batch size
(about 12 versus 5.4 min per epoch). Each run records its hardware, and
training times are compared across architectures only when measured on the
same machine.
%% TODO: add the machine used for latency if it differs from the above.
%% TODO: confirm Decisions 15--17 are approved by the team and that
%% config_final.json is sealed before submission (sellar_config_final.py).
%% TODO: the search ranges below are proposed (Decision 17); confirm approval.

\begin{table}[t]
\centering
\caption{Search space of G1--G6. The first value of each list is the initial
specification (configuration 0); the other 11 configurations are drawn from
these lists.}
\label{tab:ranges}
\footnotesize
\setlength{\tabcolsep}{3pt}
\renewcommand{\arraystretch}{0.9}
\begin{tabular}{@{}p{0.2\columnwidth}p{0.72\columnwidth}@{}}
\toprule
Model & Values \\
\midrule
G1 LSTM, G2 GRU & units $\{[64,32],[32,16],[128,64]\}$; dropout $\{0.1,0.2,0.3\}$; learning rate $\{5\!\times\!10^{-4},10^{-3},2\!\times\!10^{-3}\}$ \\
G3 CNN-1D & filters $\{[32,64,128],[16,32,64],[64,128,256]\}$; kernel $\{3,5\}$; weight decay $\{10^{-4},0,10^{-3}\}$; learning rate as G1 \\
G4 autoencoder & layers $\{[32,16,8],[64,32,16],[32,16,4]\}$; L2 $\{0,10^{-5},10^{-4}\}$; learning rate as G1 \\
G5 Transformer & blocks $\{2,1,3\}$; heads 4; $d_{model}$ $\{64,32\}$; feed-forward $\{128,64\}$; dropout $\{0.1,0.2\}$; learning rate $\{5\!\times\!10^{-4},3\!\times\!10^{-4},10^{-3}\}$; weight decay $\{10^{-4},10^{-3}\}$ \\
G6 CNN-LSTM-GRU & filters $\{64,32\}$; kernel $\{3,5\}$; LSTM units $\{64,32\}$; GRU units $\{32,16\}$; dropout $\{0.2,0.1,0.3\}$; learning rate as G1 \\
\bottomrule
\end{tabular}
\end{table}

\textbf{Selection and metrics.} The configuration maximizing mean
validation PR-AUC over the three folds is selected; the F1 threshold and the
threshold reaching recall 0.90 are chosen on the last validation block (the
final models are retrained on that block, which may bias F1 and FPR but not
PR-AUC). The primary metric is PR-AUC (average precision), whose chance level
equals the prevalence. Block-mean AP is the unweighted mean of sensor-level
AP over evaluable blocks; pooled AP is computed on scores and labels
concatenated over all valid steps of a split and is reported with its
prevalence. G0's PR-AUC depends on $w$ only; $k$ sets the binary threshold
used for F1. F1, AUC-ROC and the false-positive rate (FPR) at the F1
threshold and at the recall-0.90 threshold are secondary, and the recall
actually reached on test is reported with that FPR; PA is forbidden
\cite{kim2022rigorous,sorbo2024navigating}. Because collective anomalies
dominate the injected instants, AP is also reported per anomaly type, taking
the instants of one type as positives and the normal instants as negatives
(other types excluded). Latency is measured with \texttt{perf\_counter}
after 100 warm-up calls over $\geq$1000 calls, at batch 1 and for one window
per sensor, reporting median, p95 and 95\% CI; the operational criterion is
$\leq$10\,ms per window, complemented by the PR-AUC$\times$latency Pareto
front.

\textbf{Controls against a signature shortcut.} Networks see eight
channels, G0 only flow, and Scheme B perturbs several channels, so a network
could exploit the injection signature instead of detecting flow anomalies
better. Per-type AP and an ablation retrained with occupancy and speed zeroed
address this; the ablation keeps the hour, weekday and peak-hour channels, so
information is still not equal to G0's. Class-weighted training tests
sensitivity to imbalance, not the shortcut. These analyses were fixed after a
minimal exploratory run that is not reported as a result.

\textbf{Hypotheses and statistics.} H1: at least one DL detector exceeds G0
in block-mean test PR-AUC under Scheme B by at least $\delta=0.05$ of
absolute PR-AUC (about 14\% of G0's 0.345), a margin that rules out trivial
gains and is fixed before the test block is opened (H0: none does). H2, conditional on H1: at least one detector
satisfying H1 meets the latency criterion (median of one window, batch 1,
model inference only, common CPU; p95 also reported); H0: none does. H3: the
detector ranking is not stable across schemes (Kendall's $\tau<0.7$); H0:
$\tau\geq0.7$. Each alternative states the effect to be demonstrated (an
advantage, a latency pass, a ranking change), so H3 is phrased as
instability. Validation folds serve selection and exploratory comparisons
(Friedman test on the $291\times3$ blocks, Nemenyi
\cite{demsar2006statistical}); H1 is evaluated primarily on the 291 held-out
test blocks, after configurations, seeds, aggregation rules and comparisons
are frozen, with paired tests against G0 and Holm correction
($\alpha=0.05$). Sensors and expanding folds are not independent and Holm
does not correct dependence, so significance is indicative and effect sizes
and intervals are reported. Block AP is computed per seed and averaged over
the three seeds before testing; per-seed results are also reported. Blocks
lacking either class are excluded and reported. For H3 the conditions are
ranked by block-mean PR-AUC per scheme (blocks valid under both, resampled
for a 95\% bootstrap CI of $\tau$): the ranking is stable if the whole CI is
at or above 0.7, H3 is supported if it lies below, and the result is
inconclusive otherwise; Spearman's $\rho$ is a sensitivity check. As H3
changes labels, prevalence, training examples, selected configurations and
perturbed data together, it assesses the sensitivity of the whole procedure
to the anomaly definition, not of the evaluation labels alone. Under
Scheme A the advantage over G0 is reported descriptively only, because A
shares G0's rule family.
%% TODO (team decisions before the test block is opened for the networks):
%% (a) record as a new decision that H1 is primary on the test block, as the
%% PPI still states validation folds; (b) name the paired test (e.g. Wilcoxon
%% signed-rank); (c) confirm mean over seeds (PPI says mean or median);
%% (d) confirm median vs p95 for the H2 criterion; (e) confirm the margin
%% delta = 0.05 proposed here (about 14% of G0's 0.345) or replace it, before
%% the test block is opened; (f) state the databases, search string, dates and
%% inclusion criteria of the literature review behind "in the literature we
%% reviewed".

% ---------------------------------------------------------------------
\section{Preliminary Results}

Table~\ref{tab:results} reports the evidence obtained so far (weeks 1--8)
and what remains pending. The pipeline, both labeling schemes and G0 passed
50 of 50 automated checks, including perturbation tests that alter future
values; the window loader passed 73 of 73 and the model definitions 54 of 54
checks against the initial specification.

\begin{table}[t]
\centering
\caption{Evidence obtained to date versus pending analyses. Chance level
equals prevalence.}
\label{tab:results}
\footnotesize
\setlength{\tabcolsep}{3pt}
\renewcommand{\arraystretch}{0.9}
\begin{tabular}{@{}p{0.47\columnwidth}p{0.45\columnwidth}@{}}
\toprule
Item & Result \\
\midrule
\multicolumn{2}{@{}l}{\textit{Obtained}}\\
Sensors retained after QC & 291 of 307 (16 excluded, 5.2\%) \\
Scheme A prevalence & 0.10\% (4\,873 / 4\,860\,864); 0.086--0.125\% per partition \\
Scheme B prevalence & 3.02--3.06\% per block \\
Scheme B composition (instants) & 74.7\% collective, 14.6\% point, 10.7\% contextual \\
Overlap A vs.\ B (Jaccard) & 0.0003--0.0012 \\
Blocks with both classes & A: 340/1\,164 (29.2\%); B: 100\% \\
G0 selected config.\ (A and B) & window 22, multiplier 3.0 (multiplier at grid edge; it does not affect PR-AUC) \\
G0 test, Scheme B & block-mean PR-AUC 0.345 ($\approx$11.3$\times$ chance); pooled 0.335 (11.0$\times$); F1 0.323; AUC-ROC 0.792; FPR 54.3\% at the threshold reaching recall 0.90 on validation \\
G0 test, Scheme B, per type & point 0.905; contextual 0.399; collective 0.131 \\
G0 test, Scheme A & block-mean PR-AUC 0.084 (84 of 291 blocks evaluable); pooled 0.003 (3.0$\times$); F1 0.021; FPR 65.6\% at the validation-derived recall-0.90 threshold \\
\midrule
\multicolumn{2}{@{}l}{\textit{Pending}}\\
G1--G6 search and final training & \pend \\
Friedman/Holm, Nemenyi (H1) & \pend \\
Latency, Pareto front (H2); G7 smoke test & \pend \\
Kendall $\tau$, Spearman $\rho$ (H3) & \pend \\
Channel ablation, class weights & \pend \\
\bottomrule
\end{tabular}
\end{table}

%% TODO: add the recall reached on test at the validation-derived recall-0.90
%% threshold (printed as "recall prueba" by semana6_g0_ampliado.py).
Four observations shape the remaining work. First, the proxy and injected
labels are almost disjoint (Jaccard $\leq0.0012$), so H3 is not trivially
satisfied. Second, Scheme A has a prevalence below the anticipated 0.3--2\% and only 29.2\% of
blocks with both classes, which limits the block intersection on which H3
is computed; its two aggregations also disagree for G0 (0.084 block-mean
over 84 evaluable blocks, 0.003 pooled, about three times the pooled
prevalence), so every PR-AUC is reported with its aggregation.
Third, G0 is far above chance under injection, consistent with
the competitiveness of simple detectors
\cite{taheri2026applicability,schmidl2022anomaly}, but strongly
type-dependent: it separates point anomalies almost perfectly and collective
ones poorly, and collective anomalies are three quarters of the injected
instants, so any DL gain must be read per type. Fourth, the measured
training cost (4.3--5.4 min per epoch on CPU) forced the reduced search of
Section~III, only partly mitigated by multi-seed retraining.

% ---------------------------------------------------------------------
\section{Discussion and Threats to Validity}

Prior evidence that statistical baselines can rival DL
\cite{audibert2022deep,taheri2026applicability} motivates H1, but failure to
detect a DL advantage will be reported with effect sizes and uncertainty,
without interpreting non-significance as evidence of equivalence. The
published latency of SI-DMSeq (about 55\,ms per sequence for 46 zones and
24-step windows \cite{nouri2025sidmseq}) refers to a different load, so all
methods will be measured under one configuration; the 10\,ms criterion
(about 3\,s for one window per sensor, far below the 5-min sampling
interval) is a technical criterion of this study.

\emph{External validity:} one dataset, one freeway network and about eleven
test days limit generalization; H1 concerns injected anomalies, not verified
incidents. \emph{Construct validity:} neither scheme represents verified incidents;
injected anomalies differ from real ones \cite{carmona2022ncad}, and
unlabeled real events act as pseudo-negatives \cite{goswami2023unsupervised}.
\emph{Internal validity:} Scheme A shares its rule family with G0, hence its
restricted role; all fitting is fold-local. G0 uses flow only and a 22-step
window against 12 steps and eight channels for the networks, and the
reduced one-seed search may favor architectures that converge early. The
shortcut controls were fixed after a small exploratory run.
\emph{Conclusion validity:}
expanding folds are not independent and sensors share dynamics, so the
validation-fold Friedman analysis is exploratory and a fold-3-only analysis
is planned as a robustness check. \emph{Metric validity:} point-wise PR-AUC
weighs long (collective) anomalies more and is lag-sensitive
\cite{sorbo2024navigating,liu2024elephant}; per-type results will be
reported. \emph{Implementation validity:} G7 runs in a different framework,
a declared confounder for latency; training and latency hardware will be
documented separately.

% ---------------------------------------------------------------------
\section{Conclusion}

We specified and partially executed a uniform, reproducible protocol for
comparing DL detectors with a non-parametric baseline on PeMS-D4 freeway
flow under two separately constructed labeling schemes. Preliminary outputs
show that the schemes yield largely non-overlapping positive sets, that the
baseline is well above chance under injection yet type-dependent, and that
the nominal tuning budget was infeasible, leading to a declared reduced
search. Next come the DL battery, tests for H1--H3, shortcut controls and
the latency benchmark; code (MIT license), configurations and seeds will be
released with the results.
%% TODO: add repository URL once public.

\bibliographystyle{IEEEtran}
\bibliography{references}

\end{document}
